\begin{proof}
Fix an execution satisfying A1--A4 and a returned certificate $(\sigma,\rank,\kappa)$. We use the phase-disjoint observation map $\alpha$ from the theorem; the assumptions require it to preserve the specified component and monitor states, not merely event names. A1 maps the old prefix to same-labeled edges in the monitored old closed loop. At entry, A4 supplies one linearization point at which the actual reachable old state $z$ is observed and control changes to $\Init(z)\in Q_0$. Entry contributes the linked model's boundary edge and no game event. Thus the update observation starts in the certificate domain.

Suppose the current observation is a non-goal certificate state $q$. A2 filters controls according to $\sigma(q)$ while retaining every enabled UC event. Its next reported event and successor identity therefore give an edge $q\xrightarrow{a}q'$ in the complete policy successor relation. For a command, A3 additionally supplies the required UC quiescence, finite atomic completion and exact modeled successor. Ordinary reports satisfy A2 directly. The checked WIN certificate includes every such successor, so $q'$ remains safe and in its domain, with $\rank(q')<\rank(q)$. This establishes the update-prefix simulation and rank invariant by induction, for every outcome rather than a selected favorable one.

A2 makes a next event eventually observable at every non-goal; A3 rules out a nonterminating command between observations. Consequently the execution cannot stop or stall at a non-goal. An infinite sequence of update events would give an infinite strictly descending sequence of natural ranks, which is impossible. At most $\rank(\Init(z))$ game events therefore reach a goal. This is an event bound: the assumptions provide no uniform bound on the time between observations. They also impose no fairness on which enabled event or outcome occurs.

The goal predicate supplies UC quiescence and a compatible reachable new-endpoint state $\kappa(q)$. A4 disables ordinary controls, preserves the specified component and monitor states, and completes handover at a finite linearization point. Representing this operation by the linked boundary edge adds no game event. A1 then maps every new-phase observation to the selected monitored endpoint's same-labeled continuation. Thus the complete observed execution maps into the linked model. The argument assumes the two boundary implementations have the A4 properties; it does not prove their linearizability or atomicity from implementation code.

Finally, certificate safety preserves each supplied monitor while it is active. For each actual activation history, completeness places that history in the declared scope and RS includes the monitor's accepted continuations in the history's legal residual language. Applying this inclusion to every active interval gives the claimed requirement interpretation. The choice and intended meaning of those scopes remain input obligations. Unobservable transitions, divergence, failed commands and implementations violating A1--A4 are outside the theorem.
\end{proof}
